Drawing on my experience as the American Economic Association's  data editor, I examine the current state of open science in economics as
facilitated by and related to reproducibility. I  touch on the tension between
accessibility, sharing, and preservation. The guiding theme is the accessibility of the key ingredients for scholarship: manuscripts, data, software, and the
necessary technology to combine the latter two in order to produce knowledge.
I analyze how economic research balances openness with necessary restrictions, particularly regarding administrative and confidential data.  I  argue that a large degree of openness is nevertheless present, with extensive networks that include thousands of researchers supporting collaborative science. I  argue that resource constraints, such as software licensing costs and computational resource requirements, pose similar challenges. I  illustrate concrete benefits of open science in the economics literature, using recent articles. I wrap up by discussing the state of access to scientific articles in economics. 